\section{PS/2 键盘驱动：\texttt{keyboard.c}}

\subsection{Scancode Set 1}

PS/2 键盘按下键时发送一个 \term{make code}（0x00--0x7F），
释放时发送 \term{break code}（make code | 0x80）。

\codefile{src/kernel/drivers/keyboard.c}
\begin{lstlisting}[style=cstyle]
// Scancode Set 1 基础映射（部分）
static const char scancode_map[0x40] = {
    /* 0x00 */ 0,  0, '1','2','3','4','5','6','7','8','9','0','-','=',
    /* 0x0E */ '\b', '\t',
    /* 0x10 */ 'q','w','e','r','t','y','u','i','o','p','[',']','\n', 0,
    /* 0x1E */ 'a','s','d','f','g','h','j','k','l',';','\'','`', 0, '\\',
    /* 0x2C */ 'z','x','c','v','b','n','m',',','.','/', 0,  0,  0,  ' ',
};
\end{lstlisting}

\begin{thinkbox}
为什么键盘发送的是 scancode 而非 ASCII？

因为键盘是一个"物理开关阵列"——它只知道哪个位置的键被按了或松了。
把 scancode 翻译成哪个字符取决于\textbf{键盘布局}（US、UK、DE、中文等）。
同一个 scancode 0x10 在 QWERTY 布局是 'q'，在 AZERTY 布局是 'a'。
翻译工作由操作系统（或键盘驱动）完成。
\end{thinkbox}

\subsection{Shift 状态机}

\codefile{src/kernel/drivers/keyboard.c}
\begin{lstlisting}[style=cstyle]
static volatile uint8_t g_shift_down;  // Shift 是否被按住

static void keyboard_handle_scancode(uint8_t sc) {
    // 0xE0 前缀：扩展键（方向键等），暂时忽略
    if (sc == 0xE0u) { g_e0_prefix = 1; return; }
    if (g_e0_prefix) { g_e0_prefix = 0; return; }

    uint8_t released = (sc & 0x80u) != 0;  // 最高位=1 表示松开
    uint8_t code = sc & 0x7Fu;              // 低 7 位 = make code

    // Left Shift (0x2A) 或 Right Shift (0x36)
    if (code == 0x2Au || code == 0x36u) {
        g_shift_down = released ? 0 : 1;
        return;
    }

    if (released) return;  // 只处理按下事件

    // 根据 Shift 状态选择映射表
    char ch = g_shift_down ? scancode_map_shift[code]
                           : scancode_map[code];
    if (ch) kbd_buf_put(ch);
}
\end{lstlisting}

\subsection{环形缓冲区}

\begin{figure}[H]
  \centering
  % figures/ch05/fig04-irq-handler-head-tail.tex — auto-extracted from chapter source
\begin{tikzpicture}[node distance=0cm]
  % 缓冲区数组
  \foreach \i in {0,...,9} {
    \node[memcell, minimum width=1.2cm, minimum height=0.6cm] (buf\i) at (\i*1.3, 0) {};
    \node[font=\tiny\ttfamily, above=0.1cm of buf\i] {\i};
  }
  \node[font=\small\ttfamily, above=0.8cm of buf4, text width=8cm, align=center] {g\_kbd\_buf[128]（只画前 10 个）};

  % 内容
  \node[font=\ttfamily\small] at (buf2) {'a'};
  \node[font=\ttfamily\small] at (buf3) {'b'};
  \node[font=\ttfamily\small] at (buf4) {'c'};

  % head / tail pointers
  \draw[pointer, red!60!black] (buf2.south) -- ++(0, -0.6) node[below, font=\small] {tail（消费者读取）};
  \draw[pointer, blue!60!black] (buf5.south) -- ++(0, -0.6) node[below, font=\small] {head（IRQ 写入）};
\end{tikzpicture}

  \caption{键盘环形缓冲区：IRQ handler 写 head，主循环读 tail}
\end{figure}

\codefile{src/kernel/drivers/keyboard.c}
\begin{lstlisting}[style=cstyle]
#define KBD_BUF_SIZE 128

static volatile char    g_kbd_buf[KBD_BUF_SIZE];
static volatile uint8_t g_kbd_head;  // 下一个写位置
static volatile uint8_t g_kbd_tail;  // 下一个读位置

static void kbd_buf_put(char c) {
    uint8_t next = (g_kbd_head + 1) % KBD_BUF_SIZE;
    if (next == g_kbd_tail) return;  // 满了 → 丢弃
    g_kbd_buf[g_kbd_head] = c;
    g_kbd_head = next;
}

static int kbd_buf_get(char* out) {
    if (g_kbd_tail == g_kbd_head) return 0;  // 空
    *out = g_kbd_buf[g_kbd_tail];
    g_kbd_tail = (g_kbd_tail + 1) % KBD_BUF_SIZE;
    return 1;
}
\end{lstlisting}

\begin{designbox}
为什么用环形缓冲区而非普通数组？

\begin{enumerate}
  \item \textbf{解耦生产者和消费者}：IRQ handler（生产者）在中断上下文执行，
    必须尽快返回。Shell 主循环（消费者）在非中断上下文运行，可以慢慢处理。
  \item \textbf{避免数据拷贝}：环形缓冲区通过 head/tail 指针管理读写位置，
    不需要移动数组元素。
  \item \textbf{线程安全}：单生产者 + 单消费者的环形缓冲区在不需要锁的情况下
    就是线程安全的（原子的 head/tail 更新）。
\end{enumerate}
\end{designbox}

\begin{thinkbox}
\texttt{g\_kbd\_buf} 和 \texttt{g\_kbd\_head/tail} 为什么要声明为 \texttt{volatile}？

因为它们在两个"执行流"中被访问：
\begin{itemize}
  \item IRQ handler（中断上下文）写 \texttt{g\_kbd\_head}
  \item \texttt{keyboard\_getc}（主循环）读 \texttt{g\_kbd\_head}
\end{itemize}
没有 \texttt{volatile}，编译器可能把 \texttt{g\_kbd\_head} 缓存在寄存器中，
导致主循环永远看到旧值——缓冲区"永远为空"。
\end{thinkbox}

\subsection{公开 API}

\codefile{src/kernel/drivers/keyboard.c}
\begin{lstlisting}[style=cstyle]
// 阻塞等待一个字符（hlt 直到 IRQ 把字符放入缓冲区）
char keyboard_getc(void) {
    char c;
    while (!keyboard_try_getc(&c)) {
        __asm__ volatile("hlt");  // 等下一个中断
    }
    return c;
}

// 非阻塞尝试读取（0=无字符，1=读到一个）
int keyboard_try_getc(char* out) {
    return kbd_buf_get(out);
}
\end{lstlisting}

\begin{whybox}
为什么用 \texttt{hlt} 而非空循环？

\texttt{hlt} 让 CPU 进入低功耗状态直到下一个中断。空循环（\texttt{while(1)\{\}}）
让 CPU 100\% 占用，白白浪费电力和 QEMU 宿主机 CPU 资源。
\end{whybox}

\subsection{\texttt{keyboard\_init}：注册 + 开 IRQ}

\codefile{src/kernel/drivers/keyboard.c}
\begin{lstlisting}[style=cstyle]
void keyboard_init(void) {
    irq_install_handler(1, keyboard_irq1_handler);
    pic_clear_mask(1);  // 允许 IRQ1 通过 PIC
}
\end{lstlisting}

注意此时 EFLAGS.IF 仍然是 0（中断关闭）。IRQ handler 不会触发，
直到 \texttt{kmain} 末尾执行 \texttt{\_\_asm\_\_ volatile("sti")}。

% ============================================================================

